% ============================================================
%  Setmana 21 — Unitat 2 · Determinants
%  Exercicis: 79, 80, 83, 87
% ============================================================

\setmana{21}{Unitat 2 \textperiodcentered{} Determinants}{7 de febrer de 2027}%
        {79, 80, 83 i 87}

\begin{errada}{Exercici 80 b)}
  \ondiu{$\begin{vmatrix}3&-5&1\\4&1&-2\\7&-4&-2\end{vmatrix}=0$}
  \hadedir{$\begin{vmatrix}3&-5&1\\4&1&-2\\7&-4&-1\end{vmatrix}=0$}
  \motiu{La fila que s'ha afegit a la matriu és $(7,-4,-1)$, la suma de les
    altres dues, i per això el determinant és $0$. Amb $-2$ al final, el
    determinant valdria $-23$.}
\end{errada}

%\senseerrades{79, 83 i 87}

% ---------- Observacions (només surten amb \obstrue) ----------
\ifobs
\bigskip
\noindent
Una observació, que no és cap errada:\par\nopagebreak

\begin{observacio}{Exercici 79 d)}
  El rang 2 és correcte, però la justificació és incompleta. Dels dos determinants d'ordre 3 que es calculen, només el segon és un menor orlat del menor $\begin{vmatrix}1&3\\0&5\end{vmatrix}$ que s'agafa al final. N'hi falta l'altre, $\begin{vmatrix}2&1&3\\-2&0&5\\2&3&19\end{vmatrix}$, que també val $0$.
\end{observacio}
\fi
